% BHU PYQ -- Professional Exam Template
\documentclass[11pt,a4paper]{article}

\usepackage[a4paper,top=2.5cm,bottom=2.5cm,left=2.8cm,right=2.8cm,headheight=14pt]{geometry}
\usepackage{amsmath,amssymb}
\usepackage[T1]{fontenc}
\usepackage[utf8]{inputenc}
\usepackage{lmodern}
\usepackage{microtype}
\usepackage{fancyhdr}
\usepackage{enumitem}
\usepackage{hyperref}
\usepackage{lastpage}
\usepackage{parskip}
\usepackage{tikz}
\usetikzlibrary{arrows.meta,shapes,positioning}

\hypersetup{colorlinks=true,urlcolor=black,linkcolor=black,
  pdftitle={B.Sc. (Hons.) Zoology Sem-VI ZOB-601 -- BHU PYQ}}

\pagestyle{fancy}
\fancyhf{}
\renewcommand{\headrulewidth}{0.4pt}
\renewcommand{\footrulewidth}{0.4pt}
\fancyhead[L]{\small\textit{Banaras Hindu University}}
\fancyhead[R]{\small\textit{Zoology Paper: ZOB-601}}
\fancyfoot[C]{\small Page \thepage\ of \pageref{LastPage}}

\newlist{parts}{enumerate}{3}
\setlist[parts,1]{label=(\a*),leftmargin=*,itemsep=4pt,topsep=3pt}
\setlist[parts,2]{label=(\roman*),leftmargin=*,itemsep=2pt,topsep=2pt}
\setlist[parts,3]{label=(\arabic*),leftmargin=*,itemsep=2pt,topsep=2pt}
\newcommand{\pts}[1]{\hfill\textit{[#1]}}

\begin{document}

\begin{center}
  {\large\bfseries BANARAS HINDU UNIVERSITY}\\[2pt]
  {\normalsize B.Sc. (Hons.) Semester VI Examination, 2013-14}\\[6pt]
  \rule{0.8\linewidth}{0.5pt}\\[6pt]
  {\large\bfseries Zoology}\\[4pt]
  {\large\bfseries Paper: ZOB-601 --- Functional Anatomy and Economic Importance of Chordata}\\[4pt]
  {\normalsize Time: Three hours \hfill Full Marks: 70}
\end{center}

\rule{\linewidth}{0.4pt}

\smallskip
\noindent\textit{Note: Answer six questions including Question No. 1 which is compulsory. Draw diagrams wherever necessary. The figures in the right-hand margin indicate marks.}

\bigskip

\medskip\noindent\textbf{Question 1.} Answer the following parts:
\begin{parts}
  \item State whether the following statements are True or False, giving reasons in 2--3 sentences:\pts{$4 \times 2 = 8$}
  \begin{parts}
    \item Ostracoderms were primitive jawed vertebrates.
    \item Heterocoelous vertebra possesses convexities on both side.
    \item Osteoblast cells form dentine layer of the tooth.
    \item Spiral valve in the intestine is found in shark fishes.
  \end{parts}
  
  \item Match the terms of Column-A with Column-B:\pts{$6 \times \frac{1}{2} = 3$}
  \begin{center}
    \begin{tabular}{lp{1cm}l}
      \textbf{Column-A} & & \textbf{Column-B} \\
      (i) Pallium & & (1) Stomach \\
      (ii) Autostylic & & (2) Cerebral hemisphere \\
      (iii) Omasum & & (3) Vertebra \\
      (iv) Ctenoid & & (4) Rheoreceptors \\
      (v) Arcualia & & (5) Jaw \\
      (vi) Neuromasts & & (6) Ceruminous glands \\
       & & (7) Scale \\
    \end{tabular}
  \end{center}

  \item Differentiate between the following:\pts{$3 \times 2 = 6$}
  \begin{parts}
    \item Mullerian and Wolffian ducts
    \item Afferent and Efferent branchial arteries
    \item Nails and Hoofs.
  \end{parts}

  \item Define the following:\pts{$3 \times 1 = 3$}
  \begin{parts}
    \item Antivenin
    \item Dipnoi
    \item Holobranch.
  \end{parts}
\end{parts}

\medskip\noindent\textbf{Question 2.} Describe the different types of jaw suspensorium in vertebrates.\pts{10}

\medskip\noindent\textbf{Question 3.} Explain the following:\pts{$2 \times 5 = 10$}
\begin{parts}
  \item Types of uterus
  \item Poultry.
\end{parts}

\medskip\noindent\textbf{Question 4.} Give an account of various types of feathers found in birds.\pts{10}

\medskip\noindent\textbf{Question 5.} Discuss the changes that have occurred in cerebral hemisphere from fishes to mammals.\pts{10}

\medskip\noindent\textbf{Question 6.} Explain the following:\pts{$2 \times 5 = 10$}
\begin{parts}
  \item Dentition
  \item Urinary bladder.
\end{parts}

\medskip\noindent\textbf{Question 7.} Give an illustrated account of morphological and anatomical changes that have occurred during the course of vertebrate origin.\pts{10}

\medskip\noindent\textbf{Question 8.} Write notes on the following:\pts{$2 \times 5 = 10$}
\begin{parts}
  \item Accessory respiratory organs in fishes
  \item Role of amphibians in biological control.
\end{parts}

\vfill
\rule{\linewidth}{0.4pt}
\begin{center}\small
\href{https://bhu-exam.vercel.app/}{Click Here}\\[4pt]\href{https://me-aryan.vercel.app/}{Click Here}\\[4pt]\href{https://quashan-me.vercel.app/}{Click Here}
\end{center}

\end{document}
